\ArithmeticHead{353}{R}{{\scriptsize ZUR THEORIE DER REIHE $F(\alpha,\beta,\gamma,x)$.\quad III.}}
\begin{center}\textbf{BEMERKUNGEN.}\end{center}
\noindent\hspace*{\parindent}Neben der lateinischen Bearbeitung der \emph{Disquisitiones circa seriem}, über die \textsc{Schering} in den \textls[120]{Bemerkungen} zu dem aus dem Nachlaß herausgegebenen zweiten Teile dieser Abhandlung berichtet (Werke III, S. 230), verdient die vorstehende deutsche Bearbeitung derselben Abhandlung besondere Beachtung.
Sie besteht aus zwei Teilen, die aber kurz hintereinander niedergeschrieben sein dürften, der erste umfaßt die artt. [1.]—[7.], der zweite die artt. [8.] und [9.], beide werden wohl auf das Jahr 1809 anzusetzen sein.
Gegenüber den \emph{Disquisitiones} (Werke III, S. 132) fällt zunächst auf, daß die Beziehungen zwischen verwandten Reihen (siehe die Gleichungen (2) bis (29), oben S. 339—342) hier immer in der besonderen Form gegeben werden, daß wenn $u=F(\alpha,\beta,\gamma,x)$, $v$ eine verwandte Reihe, $u',v'$ die Derivierten bedeuten,
\[
v=Au+Bu',\quad v'=Cu+Du'
\]
\noindent gesetzt wird, wo $A,B,C,D$ gewisse rationale Funktionen von $x$ sind. Es ist also diejenige Form der Transformation, durch welche die Differentialgleichung (30), der $u$ genügt, in die (nach \textsc{Riemann}\textsuperscript{*})) \textls[120]{zu derselben Klasse} gehörige Differentialgleichung, der $v$ Genüge leistet, übergeht.

\indent Bei der Bestimmung der in der Gleichung (31) des art. [3.] auftretenden Konstanten $A$ wird auf eine spätere Stelle verwiesen.
Nun erscheint zwar im art. [7.] die analoge Gleichung [31a] für beliebige Werte der $\alpha,\beta,\gamma$ und im art. [9.] wird die Darstellung dieser allgemeinen Konstanten mit Hilfe der $\Pi$-Funktion gegeben, aber ohne Beweis.
Ein Beweis findet sich in dem nachgelassenen art. 50. der \emph{Disquisitiones}, Werke III, S. 221\quad Bemerkenswert ist im art. [3.] das Beispiel $\alpha=\tfrac12$, $\beta=\tfrac12$, $\gamma=1$, das dem arithmetisch-geometrischen Mittel entspricht; in diesem Falle ergibt sich $A=\tfrac1\pi$ und die Gleichung (31) ist nichts anderes als der \textls[120]{Schöne Lehrsatz}, oben S. 218, oder das \textls[120]{Theorema elegantissimum}, Werke VIII, S. 98.
Im art. [4.] wird die \textsc{Euler}sche Gleichung (32), vergl. oben S. 329, ähnlich wie auch in dem nachgelassenen art. 40 der \emph{Disquisitiones}, Werke III, S. 209, abgeleitet; »eine der merkwürdigsten Relationen« nennt sie \textsc{Gauss}, ohne übrigens \textsc{Euler} zu erwähnen.
Es folgen dann die Umwandlungen der Quotienten
\[
\frac{F(\alpha,\beta+1,\gamma+1)}{F(\alpha,\beta,\gamma)}\quad\text{und}\quad\frac{F(\alpha+1,\beta,\gamma)}{F(\alpha,\beta,\gamma)}
\]
\noindent in Kettenbrüche, ganz ähnlich wie im art. 12. der \emph{Disquisitiones}, Werke III, S. 134. Die Frage der
\par\medskip\noindent\rule{.25\linewidth}{.4pt}\par\smallskip
{\small\noindent *) B. \textsc{Riemann}, Gesammelte Werke, 2. Aufl., 1892, S. 380.\par}
\par\bigskip\noindent\small X1.\hfill45\normalsize
\clearpage\SourceLeaf{361}\ArithmeticHead{354}{L}{ANALYSIS.\quad NACHLASS.}
\noindent Konvergenz dieser Kettenbrüche hat \textsc{Gauss} nicht erörtert, während er die Konvergenz der Reihe $F(\alpha,\beta,\gamma,x)$ namentlich auch für den Punkt $x=1$ in der Abhandlung von 1812 genau untersucht.
Daß er aber über die formale Umformung hinaus auch die quantitative Bedeutung der Kettenbruchentwicklung erkannt hat, ergibt sich schon aus dem oben S. 237 abgedruckten Briefe an \textsc{Bessel}, in dem die Brauchbarkeit der dort betrachteten besonderen Kettenbrüche für die numerische Berechnung maßgebend ist. — Wie L. W. \textsc{Thomé}\textsuperscript{*}) gezeigt hat, konvergiert der Kettenbruch, der dem Quotienten
\[
\frac{F(\alpha,\beta+1,\gamma+1)}{F(\alpha,\beta,\gamma)}
\]
\noindent entspricht, in der ganzen komplexen $x$-Ebene mit Ausnahme der reellen Werte, die nicht kleiner sind als 1, und stellt dort den Quotienten jener Reihen bezw. den ihrer analytischen Fortsetzungen innerhalb der durch den Querschnitt $(1,+\infty)$ zerschnittenen Ebene dar.

\indent Die artt. [5.], [6.] geben die Begründung der in dem Briefe an \textsc{Bessel}, oben S. 237, auseinandergesetzten Entwicklungen, wobei wieder das agM. eine Rolle spielt, das diesmal auch ausdrücklich genannt wird.
Die am Schluß von art. [7.] später hinzugeschriebenen Formeln enthalten schon die erst im art. [8.] eingeführte $\Pi$-Funktion.
In den Gleichungen I., II. des art. [8.], die mit den Gln. [100], [101] Werke III, S. 225 gleichbedeutend sind, entspricht die angewandte Transformation von $x^2$ in $\frac{4x}{(1+x)^2}$ dem Übergang von $x=\frac ba$ zu $x_1=\frac{b_1}{a_1}$, wo
\[
a_1=\frac{a+b}{2},\quad b_1=\sqrt{ab};
\]
\noindent vergl. auch in der Abhandlung von 1812, Werke III, S. 129 die \textls[120]{dritte} Form (»tertio fit«) der dort auftretenden Koeffizienten.
In der Gleichung IV. kann $i$ einen ganz beliebigen Wert haben.
Zu V. vergleiche man Werke III, S. 227, Gleich. [106] und S. 228, Gleich. [107].
Für den am Schluß von art. [9.] behandelten Ausnahmefall $\gamma=1$, der auch beim agM. auftritt, vergl. den nachgelassenen art. 45 der \emph{Disquisitiones}, Werke III, S. 214, für art. [10.] die artt. [18.] und ff. der \emph{Disquisitiones}, ebenda S. 144.
\par\hfill\textsc{Schlesinger.}\par
\medskip\noindent\rule{.25\linewidth}{.4pt}\par\smallskip
{\small\noindent *) \emph{Über die Kettenbruchentwicklung der Gaussschen Function} $F(\alpha,1,\gamma,x)$, \textsc{Crelles} Journal für Mathematik \textbf{66}, 1866, S. 322 und \emph{Über die Kettenbruchentwicklung des Gaussschen Quotienten} $F(\alpha,\beta+1,\gamma+1,x):F(\alpha,\beta,\gamma,x)$, ebenda \textbf{67}, 1867, S. 299; siehe auch \textsc{Riemanns} aus dem Jahre 1863 stammendes Fragment \emph{Sullo svolgimento del quoziente di due serie ipergeometriche in frazione continua infinita}, in der Bearbeitung von H. A. \textsc{Schwarz} zuerst veröffentlicht 1876 in der 1. Auflage von \textsc{Riemanns} Gesammelten Werken, 2. Auflage, 1892, S. 424.\par}
\begin{center}\rule{.25\textwidth}{.4pt}\end{center}
\clearpage
